\section{Conclusion}
\label{sec:conclusion}

We presented Kimodo-Policy, a humanoid vision–language–action framework that transfers pretrained text-to-motion priors to visually grounded, closed-loop whole-body control. Its root–body–hand architecture combines dedicated motion generators with explicit cross-stage conditioning, accommodating distinct control requirements while maintaining coordinated action generation. Trainable visual branches ground the pretrained motion backbone in the current scene, while motion-history conditioning enables continual replanning from updated observations. Experiments on HumanoidArena, SIMPLE, and eight real-world tasks demonstrate strong performance across diverse locomotion and manipulation scenarios, with only 92M trainable parameters out of 678M total parameters. Further analyses show consistent benefits from stronger T2M backbones, larger VLA pretraining datasets, and increased policy capacity within the evaluated settings. Together, these results establish pretrained motion generators as an effective foundation for humanoid VLA learning, allowing policy adaptation to build on existing coordination and temporal priors when learning task-conditioned interaction.

\textbf{Limitations.} A principal limitation of the current framework is its dependence on a specific low-level motion tracker. Kimodo-Policy is designed and evaluated with SONIC, and its motion-reference interface and deployment pipeline are tailored to this tracker; other tracking controllers are not currently supported. Consequently, the reported results establish the effectiveness of the policy within the SONIC-based control stack, while transfer across trackers remains unverified. Although separating high-level motion generation from low-level execution provides a modular organization, this separation alone does not ensure compatibility with controllers that differ in reference parameterization, tracking behavior, or executable motion range. Extending the framework would therefore require adapting the policy–tracker interface and validating closed-loop performance under the target controller. Developing a more general interface between pretrained motion generation and low-level whole-body control remains an important direction for broadening its applicability.

\textbf{Future Work.} Looking ahead, Kimodo-Policy can be viewed as a motion-prior-informed action expert within a larger multimodal architecture. Integrating it with vision–language models and video generation models, potentially through MoT-style architectures, could combine semantic understanding and visual dynamics modeling with structured whole-body action generation. A complementary direction is to strengthen the underlying motion prior itself. Text–motion supervision can draw on human motion data retargeted to robot embodiments without requiring every sequence to be collected through robot execution with synchronized visual observations, offering a scalable route to motion pretraining. The gains observed with stronger T2M backbones motivate investigating whether broader motion coverage and more capable generators yield further improvements in downstream control. Beyond training robot-space motion models on retargeted data, a longer-term direction is to transfer pretrained human T2M backbones directly into humanoid VLA policies, drawing on the breadth of human motion resources and advances in human motion generation. Such transfer would require aligning motion representations across embodiments and adapting human kinematic priors to robot-specific actuation, balance, and contact constraints. These directions could expand the framework from a standalone policy into a reusable motion-generation component for more capable humanoid systems.